We report results for correlation existence (RQ1, P1) and clustering quality (RQ2, P2). Scale invariance (RQ3, P3) and intervention validation (RQ4, P4) were not evaluated due to clustering gate failure. All reported p-values use Bonferroni correction ($\alpha = 0.0033$ for three pairwise tests).

\subsection{Correlation Existence (RQ1, P1): VALIDATED}

\textbf{Finding:} All three benchmark pairs exhibit near-perfect positive correlations ($r > 0.99$, $p < 2 \times 10^{-17}$), far exceeding the hypothesized medium effect size ($r > 0.3$).

\begin{table}[h]
\centering
\begin{tabular}{lccc}
\toprule
\textbf{Benchmark Pair} & \textbf{Spearman r} & \textbf{p-value (Bonferroni)} & \textbf{Significance} \\
\midrule
TrustfulQA $\leftrightarrow$ AdvBench & 0.998 & $1.11 \times 10^{-23}$ & \checkmark \\
TrustfulQA $\leftrightarrow$ BOLD & 0.993 & $1.35 \times 10^{-17}$ & \checkmark \\
AdvBench $\leftrightarrow$ BOLD & 0.996 & $9.96 \times 10^{-20}$ & \checkmark \\
\bottomrule
\end{tabular}
\caption{Pairwise correlations between trustworthiness benchmarks}
\end{table}

\textbf{Interpretation:} Correlations $3 \times$ stronger than hypothesized threshold ($r > 0.99$ vs $r > 0.3$) indicate trustworthiness failures are not moderately coupled but near-redundant---models failing TrustfulQA (reliability) almost certainly fail AdvBench (robustness) and BOLD (fairness) at nearly identical rates. This challenges the independent dimensions assumption underlying current evaluation practice.

\textbf{Gate Metrics:}

\begin{table}[h]
\centering
\begin{tabular}{lccc}
\toprule
\textbf{Metric} & \textbf{Target} & \textbf{Actual} & \textbf{Status} \\
\midrule
Significant Pairs & $\geq 1$ (PoC), $\geq 2$ (full) & 3/3 & \checkmark\checkmark \\
Mean Effect Size & $>0.3$ & 0.996 & \checkmark\checkmark \\
Min Effect Size & $>0.3$ & 0.993 & \checkmark\checkmark \\
Bonferroni Pass Rate & $>0.67$ & 1.0 & \checkmark\checkmark \\
\bottomrule
\end{tabular}
\caption{Gate metrics for correlation analysis}
\end{table}

All gate metrics exceeded both proof-of-concept and full validation criteria. The 100\% pass rate (3/3 significant pairs) and mean $r = 0.996$ suggest dimensional near-redundancy rather than moderate coupling from shared root causes.

\textbf{Stratified Analysis (Controlling for Model Size):}

To rule out model size confounds, we computed correlations separately within size strata:

\begin{table}[h]
\centering
\small
\begin{tabular}{lcccc}
\toprule
\textbf{Stratum} & \textbf{n} & \textbf{TrustfulQA $\leftrightarrow$ AdvBench} & \textbf{TrustfulQA $\leftrightarrow$ BOLD} & \textbf{AdvBench $\leftrightarrow$ BOLD} \\
\midrule
Small ($<$1B) & 6 & $r=0.994$, $p=2.64 \times 10^{-5}$ & $r=0.989$, $p=8.12 \times 10^{-5}$ & $r=0.986$, $p=1.45 \times 10^{-4}$ \\
Medium (1-10B) & 9 & $r=0.998$, $p=1.47 \times 10^{-9}$ & $r=0.995$, $p=1.31 \times 10^{-8}$ & $r=0.997$, $p=4.37 \times 10^{-9}$ \\
Large ($>$10B) & 5 & $r=0.997$, $p=2.39 \times 10^{-4}$ & $r=0.991$, $p=1.01 \times 10^{-3}$ & $r=0.994$, $p=4.58 \times 10^{-4}$ \\
\bottomrule
\end{tabular}
\caption{Stratified correlations within model size groups}
\end{table}

\textbf{Key observation:} Correlations remain $r > 0.98$ within all strata, indicating coupling persists even when comparing models of similar size. This rules out the alternative explanation that correlations merely reflect "larger models are uniformly better across dimensions." The tight coupling generalizes across model scales (partial evidence for P3 despite blocked Mantel test).

\subsection{Clustering Quality (RQ2, P2): REFUTED}

\textbf{Finding:} Hierarchical clustering produces perfectly stable clusters (100\% bootstrap consistency) but poor separation (silhouette = 0.274 $<$ 0.5 threshold).

\begin{table}[h]
\centering
\begin{tabular}{lccc}
\toprule
\textbf{Metric} & \textbf{Value} & \textbf{Threshold} & \textbf{Status} \\
\midrule
Optimal k & 2 & 2-5 & \checkmark \\
Silhouette score & 0.274 & $>0.5$ & $\times$ \\
Bootstrap consistency & 100.0\% & $\geq 80\%$ & \checkmark \\
Cophenetic correlation & 0.693 & $>0.7$ & $\times$ (borderline) \\
\bottomrule
\end{tabular}
\caption{Clustering quality metrics}
\end{table}

\textbf{Cluster Assignments ($k=2$):}
\begin{itemize}
\item \textbf{Cluster 1:} TrustfulQA, AdvBench
\item \textbf{Cluster 2:} BOLD
\end{itemize}

\textbf{Bootstrap Consistency Matrix} (\% co-occurrence across 1000 iterations):

\begin{table}[h]
\centering
\begin{tabular}{lccc}
\toprule
 & \textbf{TrustfulQA} & \textbf{AdvBench} & \textbf{BOLD} \\
\midrule
\textbf{TrustfulQA} & 100\% & 100\% & 0\% \\
\textbf{AdvBench} & 100\% & 100\% & 0\% \\
\textbf{BOLD} & 0\% & 0\% & 100\% \\
\bottomrule
\end{tabular}
\caption{Bootstrap consistency matrix}
\end{table}

\textbf{Interpretation:} Clusters are perfectly stable (TrustfulQA+AdvBench always cluster together, BOLD always separate across all 1000 bootstrap iterations) but not well-separated in silhouette metric (0.274 $\ll$ 0.5). This divergence reveals that bootstrap consistency (stability) and silhouette score (separation) measure orthogonal properties:
\begin{itemize}
\item \textbf{Stability:} Do clusters remain consistent across resampling? (100\% = yes)
\item \textbf{Separation:} Are clusters far apart in distance space? (0.274 = no)
\end{itemize}

With $r > 0.99$ correlations, all benchmarks are nearly collinear, yielding minimal distance variation (0.003-0.007 range). Clusters are distinguishable (stable) but distances are so small that silhouette---which penalizes small inter-cluster gaps---scores them as poorly separated.

\textbf{Methodological Constraint:} $k \geq 3$ silhouette evaluation not computable with $n=3$ benchmarks (sklearn requires $k <$ n\_samples). Testing the original hypothesis prediction of 2-5 distinct failure modes requires expanding to $\geq 5$ benchmarks.

\textbf{Gate Decision:} $\times$ FAIL (MUST\_WORK gate not met)

Clustering failed the critical separation threshold (silhouette $< 0.5$), blocking downstream hypotheses:
\begin{itemize}
\item \textbf{H-M4 (scale invariance):} Cannot test whether unstable clusters persist across strata.
\item \textbf{Intervention validation (P4):} Cannot target specific failure modes without validated taxonomy.
\end{itemize}

\textbf{Figure 1:} Dendrogram shows hierarchical clustering structure (TrustfulQA+AdvBench merge at $d=0.002$, then merge with BOLD at $d=0.007$). Despite clear visual separation, silhouette metric flags poor quality due to small absolute distances.

\textbf{Figure 2:} Silhouette plot for $k=2$ shows both clusters have positive but low silhouette coefficients (Cluster 1: $s=0.28$, Cluster 2: $s=0.24$), confirming borderline separation.

\textbf{Figure 3:} Consistency matrix heatmap visualizes 100\% bootstrap stability (diagonal blocks perfectly consistent).

\subsection{Competing Explanations for $r > 0.99$ Coupling}

The near-perfect correlations admit two interpretations:

\subsubsection{Hypothesis 1: Unified Capability}

\textbf{Claim:} All three benchmarks measure the same underlying construct (general model quality or unified trustworthiness) rather than orthogonal dimensions.

\textbf{Supporting Evidence:}
\begin{itemize}
\item $r > 0.99$ across all pairs (near-redundancy)
\item $k=2$ clustering (TrustfulQA+AdvBench vs BOLD) lacks semantic interpretation---no clear mapping to epistemic uncertainty vs distribution shift vs bias amplification
\item Stratified analysis shows slight variation ($r=0.986$-0.998) but not enough to distinguish dimensions
\item Alignment with prior observations that RLHF improves multiple dimensions simultaneously
\end{itemize}

\textbf{Contradicting Evidence:}
\begin{itemize}
\item Different benchmark formats (TrustfulQA: MC questions, AdvBench: attack success rate, BOLD: bias metrics) yet correlations persist
\item Correlation does not prove causal unity---shared confounds (e.g., all benchmarks sensitive to training data diversity) could produce $r > 0.99$ without unified construct
\end{itemize}

\textbf{Discriminating Prediction:} 10-benchmark analysis shows $r > 0.95$ for $>80\%$ of benchmark pairs AND factor analysis reveals single dominant component explaining $>90\%$ variance.

\subsubsection{Hypothesis 2: Insufficient Resolution}

\textbf{Claim:} Distinct failure modes exist but 3 benchmarks cannot resolve them---broader measurement (5-10 benchmarks) would break correlations into separable clusters.

\textbf{Supporting Evidence:}
\begin{itemize}
\item 100\% bootstrap consistency shows clusters are stable (reproducible structure)
\item HELM uses 50+ benchmarks---our 3-benchmark subset may capture shared variance but miss dimension-specific patterns
\item $k=3-5$ evaluation blocked by $n=3$ sample size constraint
\end{itemize}

\textbf{Contradicting Evidence:}
\begin{itemize}
\item $r > 0.99$ is near-perfect---even 10 benchmarks unlikely to break this into distinct clusters if underlying factor is unified
\end{itemize}

\textbf{Discriminating Prediction:} 10-benchmark analysis shows mean $r < 0.85$ across pairs with $\geq 30\%$ of pairs exhibiting $r < 0.7$ AND clustering yields well-separated groups (silhouette $> 0.5$).

\subsection{Summary of Validated and Refuted Claims}

\begin{table}[h]
\centering
\small
\begin{tabular}{lcccc}
\toprule
\textbf{Prediction} & \textbf{Hypothesized} & \textbf{Observed} & \textbf{Status} & \textbf{Confidence Adjustment} \\
\midrule
\textbf{P1:} Correlations exist & $r > 0.3$, $p < 0.01$ & $r > 0.99$, $p < 2 \times 10^{-17}$ & \checkmark SUPPORTED & +0.40 (exceeded) \\
\textbf{P2:} Distinct clusters & silhouette $> 0.5$ & silhouette = 0.274 & $\times$ REFUTED & -0.45 (failed) \\
\textbf{P3:} Scale invariance & Mantel $r > 0.7$ & NOT EVALUATED & $\triangle$ UNTESTED & -0.40 (blocked) \\
\textbf{P4:} Intervention & $d > 0.5$ & NOT EVALUATED & $\triangle$ UNTESTED & 0 (depends on P2) \\
\bottomrule
\end{tabular}
\caption{Summary of hypothesis predictions and outcomes}
\end{table}

\textbf{Overall Hypothesis Status:} CORRELATION\_VALIDATED\_TAXONOMY\_REFUTED (1/4 predictions validated, 1/4 refuted, 2/4 untested due to methodology constraints rather than evidence refuting them)

\textbf{Revised Confidence:} 0.85 $\rightarrow$ 0.40
\begin{itemize}
\item \textbf{+0.40:} P1 far exceeded threshold ($r > 0.99$ instead of $r > 0.3$)
\item \textbf{-0.45:} P2 failed critical gate (silhouette $< 0.5$)
\item \textbf{-0.40:} P3 blocked by methodology (not refuted by data)
\end{itemize}

The correlation finding is robust ($r > 0.99$, $p < 2 \times 10^{-17}$, generalizes across strata), but the failure mode taxonomy is not validated (clustering quality insufficient). This shifts the contribution from "distinct failure modes discovered" to "tight coupling revealed, taxonomy validation requires broader measurement."
